%% Non-redundancy of the Casas--Alvero conditions
%%
%% Generated by Claude Opus 5 (Anthropic) on 2026-09-16 (Spanish) and translated by the
%% same model on 2026-09-17.  Published unverified in AI-Generated Mathematical
%% Explorations, exploration 002.
%%
%% This file is the AI output as generated, with two editorial changes, both made
%% by the curator and listed here in full:
%%   1. this header;
%%   2. a provenance note added after \maketitle, so that the PDF carries its own
%%      status when it travels alone; the line "% TODO: authors" of the original was
%%      removed with it.
%% The mathematical text is untouched.  The original files, Spanish and English, are
%% reproduced verbatim in ai-output.md, next to this one.

\documentclass[11pt,reqno]{amsart}

\usepackage[T1]{fontenc}
\usepackage[utf8]{inputenc}
\usepackage{amsmath,amssymb,amsthm,mathtools}
\usepackage{array,booktabs}
\usepackage[margin=1.1in]{geometry}
\usepackage[colorlinks=true,linkcolor=blue,citecolor=blue,urlcolor=blue]{hyperref}

\theoremstyle{plain}
\newtheorem{theorem}{Theorem}[section]
\newtheorem{proposition}[theorem]{Proposition}
\newtheorem{lemma}[theorem]{Lemma}
\newtheorem{corollary}[theorem]{Corollary}
\newtheorem{conjecture}[theorem]{Conjecture}
\theoremstyle{definition}
\newtheorem{definition}[theorem]{Definition}
\newtheorem{remark}[theorem]{Remark}

\DeclareMathOperator{\Res}{Res}
\newcommand{\KK}{\overline{K}}

\title[Non-redundancy of the Casas--Alvero conditions]
      {Non-redundancy of the Casas--Alvero conditions:\\
       an elementary lemma and its consequences}
\author{}
\date{September 17, 2026}

\begin{document}

\begin{abstract}
Let $f=x^d+a_1x^{d-1}+\dots+a_{d-1}x$ be generic and let
$R_i=\Res(f,H_i(f))$ be the resultant of $f$ with its $i$-th Hasse derivative.
Schaub and Spivakovsky proved that $R_i\notin\sqrt{(R_j:j\neq i)}$ for
$i\in\{d-3,d-2,d-1\}$. We show that this holds for \emph{every}
$i\in\{1,\dots,d-1\}$, with an elementary proof based on explicit witnesses: on
the lines $t\mapsto x^d+t\,x^i$ (and $t\mapsto x^d-t\,x^{d-1}$) all the
resultants $R_j$ with $j\neq i$ vanish and
$R_i=\pm\big(\binom{d}{i}-1\big)^{d-i}t^d$. We also discuss the consequences:
none of the $d-1$ Casas--Alvero conditions is redundant; the conjecture is
equivalent to $R_i$ not vanishing identically on \emph{any} component of
$V(R_j:j\neq i)$, and our lemma establishes this for \emph{one}; and every
membership certificate $a_k^N\in(R_1,\dots,R_{d-1})$ is strongly constrained
along those lines.
\end{abstract}

\maketitle

\begin{center}\small
\fbox{\begin{minipage}{0.92\textwidth}
\textbf{Provenance.} This note was generated by an AI system (Claude Opus 5,
Anthropic) on 16 September 2026 and translated into English by the same
system the following day, in the context of research activities of
Santiago Laplagne at the Instituto de C\'alculo, Universidad de Buenos Aires. It
is published \textbf{unverified}: no human author claims these mathematical
results, and the arguments have not been independently checked. Current status,
original AI output, and how to report an error or a verification:
\url{https://github.com/slap/ai-mathematical-explorations}
\end{minipage}}
\end{center}

\section{Setting}

We follow the notation of Schaub--Spivakovsky \cite{SS}. Let $K$ be a field of
characteristic $0$, let $d\ge3$, and let
\[
  f=x^d+a_1x^{d-1}+\dots+a_{d-1}x\ \in\ K[a_1,\dots,a_{d-1}][x],
\]
that is, with $a_0=1$ and $a_d=0$ (after a translation one may always assume that
$0$ is a root of $f$). For $1\le i\le d-1$ let
\[
  H_i(f)=\sum_{k=0}^{d-i}\binom{d-k}{i}a_k\,x^{d-k-i}
\]
be the $i$-th Hasse derivative, $H_i(f)=f^{(i)}/i!$, and set
\[
  R_i:=\Res\big(f,H_i(f)\big)\in K[a_1,\dots,a_{d-1}],\qquad
  J:=(R_1,\dots,R_{d-1}),\qquad J_i:=(R_j:\ j\neq i).
\]
A polynomial $f$ has a non-constant common factor with $H_i(f)$ if and only if
$R_i(f)=0$.

\begin{conjecture}[Casas--Alvero]\label{conj:CA}
If $f$ has a non-constant common factor with every $H_i(f)$, $1\le i\le d-1$,
then $f=x^d$. Equivalently, $V(J)=\{0\}$; equivalently, $a_k^N\in J$ for every
$k$ and some $N$.
\end{conjecture}

Since the $R_j$ have integer coefficients, the validity of the last formulation
does not depend on $K$ (in characteristic $0$), and for that reason everything
below is over $\mathbb{Q}$. Schaub and Spivakovsky also observe
\cite[Rmk.~4]{SS} that the conjecture is equivalent to $R_1,\dots,R_{d-1}$ being
a regular sequence in $K[a_1,\dots,a_{d-1}]$, in any order.

\subsection*{Quasi-homogeneity}
With the weights $\deg a_k=k$, each $R_j$ is quasi-homogeneous: the substitution
$a_k\mapsto\lambda^ka_k$ corresponds to $f(x)\mapsto\lambda^df(x/\lambda)$, which
rescales the roots of $f$ and of $H_j(f)$ by $\lambda$. Consequently $V(J)$ and
$V(J_i)$ are cones, stable under the action of $\KK^{\,*}$.

\section{The result of Schaub--Spivakovsky}

\begin{theorem}[{\cite[Theorem~5]{SS}}]\label{thm:SS}
If $i\in\{d-3,d-2,d-1\}$ then
$R_i\notin\sqrt{(R_1,\dots,\check R_i,\dots,R_{d-1})}$.
\end{theorem}

This is a partial result ``in the direction'' of the conjecture: if the
conjecture holds then $R_i\notin\sqrt{J_i}$ for \emph{every} $i$, because if
$R_i$ were in the radical of the others we would have $V(J_i)=V(J)$, and the
height of $J$ would be at most $d-2$.

The proof in \cite{SS} is indirect. Denying $R_i\in\sqrt{J_i}$ amounts to
exhibiting $f$ with $R_j(f)=0$ for $j\neq i$ and $R_i(f)\neq0$. Draisma and de
Jong \cite{DdJ} construct, for each $i$, real polynomials satisfying all the
conditions $j\neq i$ which are not powers of a linear polynomial (\emph{almost
counterexamples of level $i$}). But ``not being a power of a linear polynomial''
does not imply $R_i(f)\neq0$ without already knowing the conjecture. To rule out
those polynomials being genuine counterexamples, \cite{SS} combines arguments on
the location of real roots (Rolle) with the fact that a counterexample has at
least five distinct roots \cite{CLO}; those arguments only close for the three
highest indices.

\section{The lemma: all indices}

The idea is to use \emph{degenerate} witnesses: a root $0$ of multiplicity $i$ is
shared by all the derivatives except the $i$-th, and checking $R_i\neq0$ becomes
a direct computation.

\begin{proposition}[values on the witness lines]\label{prop:lines}
For $t\in K$ let
\[
  f_{i,t}:=x^d+t\,x^i=x^i\big(x^{d-i}+t\big)\quad(1\le i\le d-2),\qquad
  f_{d-1,t}:=x^d-t\,x^{d-1}=x^{d-1}(x-t).
\]
Then, for every $i\in\{1,\dots,d-1\}$,
\[
  R_j(f_{i,t})=0\ \ (j\neq i),\qquad
  R_i(f_{i,t})=\pm\Big(\binom{d}{i}-1\Big)^{d-i}t^{d}.
\]
\end{proposition}

\begin{proof}
We use Poisson's formula
$\Res(f,g)=\pm\,\mathrm{lc}(g)^{\deg f}\prod_{g(\eta)=0}f(\eta)$.

\emph{Case $1\le i\le d-2$.} Here
$H_j(f_{i,t})=\binom dj x^{d-j}+t\binom ij x^{i-j}$.
\begin{itemize}
\item If $j<i$, both summands vanish at $0$, and $f_{i,t}(0)=0$; hence $R_j=0$.
\item If $j>i$, then $\binom ij=0$ and $H_j=\binom dj x^{d-j}$ with $d-j\ge1$; it
  shares the root $0$, hence $R_j=0$.
\item If $j=i$, set $C=\binom di$ and $m=d-i$. Then $H_i=Cx^m+t$, whose roots
  $\eta_1,\dots,\eta_m$ satisfy $\eta^m=-t/C$ and $\prod\eta_\ell=\pm t/C$. At
  those roots
  \[
    f_{i,t}(\eta)=\eta^i(\eta^m+t)=\eta^i\,t\Big(1-\frac1C\Big),
  \]
  so that
  \[
    R_i=\pm\,C^{d}\Big(\prod_\ell\eta_\ell\Big)^{i}t^m\Big(\frac{C-1}{C}\Big)^m
       =\pm\,C^{i+m}\,\frac{t^i}{C^i}\,t^m\,\frac{(C-1)^m}{C^m}
       =\pm\,(C-1)^{d-i}\,t^d .
  \]
\end{itemize}

\emph{Case $i=d-1$.} Here
$H_j(f_{d-1,t})=\binom dj x^{d-j}-t\binom{d-1}{j}x^{d-1-j}$. If $j\le d-2$ both
exponents are $\ge1$, so $H_j$ shares the root $0$ with $f$ and $R_j=0$. If
$j=d-1$, then $H_{d-1}=dx-t$ has root $t/d$ and
$f_{d-1,t}(t/d)=(t/d)^{d-1}(t/d-t)=-(d-1)t^d/d^d$, whence
$R_{d-1}=\pm d^d\cdot(d-1)t^d/d^d=\pm(d-1)t^d$, which is the formula with
$\binom{d}{d-1}-1=d-1$.
\end{proof}

\begin{theorem}\label{thm:all}
For every $d\ge3$ and every $i\in\{1,\dots,d-1\}$,
\[
  R_i\notin\sqrt{(R_1,\dots,\check R_i,\dots,R_{d-1})}\quad\text{in }K[a_1,\dots,a_{d-1}].
\]
\end{theorem}

\begin{proof}
Suppose $R_i^N=\sum_{j\neq i}c_jR_j$ with $c_j\in K[a]$. Evaluating at the
$K$-rational point corresponding to $f_{i,1}$, Proposition~\ref{prop:lines} gives
$0\neq\big(\pm(\binom di-1)^{d-i}\big)^N=0$, since $\binom di\ge d\ge3$ for
$1\le i\le d-1$. Contradiction.
\end{proof}

\begin{remark}
The proof does not use the Nullstellensatz: only the trivial direction ``if $R_i$
lies in the radical then it vanishes at every zero of $J_i$''. It uses neither
real roots nor the conjecture. In coordinates, the witness of index $i\le d-2$ is
the point with $a_{d-i}=1$ and all other coordinates zero; the one of index $d-1$
is $a_1=-1$. In both cases the coordinate that moves is $a_{d-i}$.
\end{remark}

\begin{remark}
Theorem~\ref{thm:all} has been formally verified in Lean~4 with Mathlib, for
every $d$ and every $i$, over an arbitrary field of characteristic zero: see the
theorems \texttt{exists\_witness} and \texttt{exists\_witness\_resultant} in the
accompanying development.
\end{remark}

\section{Consequences}

\subsection{No condition is redundant}

\begin{corollary}\label{cor:nonredundant}
For each $i$, the $d-2$ conditions ``$f$ shares a root with $H_j(f)$, $j\neq i$''
do not imply the $i$-th one, not even at the level of radicals. More precisely,
the line $L_i:=\{f_{i,t}:t\in K\}$ is contained in $V(J_i)$ and
$L_i\cap V(R_i)=\{0\}$.
\end{corollary}

Note the difference with the almost counterexamples of \cite{DdJ}: the
polynomials $f_{i,1}$ not only satisfy the conditions $j\neq i$, they
\emph{verifiably violate} the $i$-th one.

\subsection{Where the difficulty lies: one component against all}

\begin{lemma}\label{lem:components}
Fix $i$. Conjecture~\ref{conj:CA} holds in degree $d$ if and only if every
irreducible component $C$ of $V_{\KK}(J_i)$ satisfies $\dim C=1$ and
$R_i|_C\not\equiv0$.
\end{lemma}

\begin{proof}
($\Leftarrow$) Each component $C$ is stable under the action of $\KK^{\,*}$ (the
group is connected) and contains $0$ (the limit $\lambda\to0$). If $\dim C=1$ and
$p\in C\smallsetminus\{0\}$, the orbit of $p$ is infinite (the weights are
positive), so its closure is an irreducible closed set of dimension $1$ contained
in $C$, hence equal to $C$. Since $R_i$ is quasi-homogeneous,
$R_i(\lambda\cdot p)=\lambda^wR_i(p)$; if $R_i(p)=0$ for some $p\neq0$ in $C$
then $R_i$ vanishes on the whole orbit, hence on $C$, contrary to hypothesis.
Therefore $C\cap V(R_i)=\{0\}$ for every component, and
$V(J)=V(J_i)\cap V(R_i)=\{0\}$.

($\Rightarrow$) Suppose $V(J)=\{0\}$. By Krull's principal ideal theorem, every
component of $V(J_i)$ (defined by $d-2$ equations in $d-1$ variables) has
dimension $\ge1$. If some $C$ had dimension $\ge2$, then $C\cap V(R_i)$ would be
nonempty (it contains $0$) of dimension $\ge1$, contained in $V(J)$, absurd. And
if $R_i\equiv0$ on some $C$, then $C\subseteq V(J)$, also absurd.
\end{proof}

Theorem~\ref{thm:all} says that the component of $V(J_i)$ containing $L_i$
satisfies $R_i|_C\not\equiv0$. The conjecture asks for the same on \emph{every}
component (and for all of them to be curves). In particular, if the conjecture
holds then $L_i$ is an irreducible component of $V(J_i)$. The distance between
``one component'' and ``all of them'' is exactly the difficulty of
Casas--Alvero.

\subsection{Constraints on membership certificates}

The membership formulation of the conjecture asks, for each $k$, for a
certificate
\begin{equation}\label{eq:cert}
  a_k^N=\sum_{j=1}^{d-1}c_j\,R_j,\qquad c_j\in\mathbb{Q}[a_1,\dots,a_{d-1}].
\end{equation}
Let $\varphi_i:\mathbb{Q}[a]\to\mathbb{Q}[t]$ be the restriction to the line
$L_i$, that is, $a_{d-i}\mapsto t$ (or $a_1\mapsto-t$ if $i=d-1$) and all other
coordinates $\mapsto0$. Each coordinate $a_k$ moves on exactly one of the lines,
the one of index $i=d-k$.

\begin{corollary}\label{cor:certificates}
Every certificate \eqref{eq:cert} satisfies:
\begin{enumerate}
\item $N\ge d$;
\item $\varphi_{d-k}(c_{d-k})=\pm\,t^{N-d}\big/\big(\binom dk-1\big)^{k}$; in
  particular, the certificate for $a_k^N$ \emph{must use} $R_{d-k}$;
\item $\varphi_i(c_i)=0$ for every $i\neq d-k$.
\end{enumerate}
\end{corollary}

\begin{proof}
Apply $\varphi_i$ to \eqref{eq:cert}. By Proposition~\ref{prop:lines} only the
term $j=i$ survives:
$\varphi_i(a_k)^N=\varphi_i(c_i)\cdot\big(\pm(\binom di-1)^{d-i}t^d\big)$. If
$i=d-k$, the left-hand side is $\pm t^N$, which forces $N\ge d$ and the value of
$\varphi_{d-k}(c_{d-k})$, using $\binom d{d-k}=\binom dk$. If $i\neq d-k$, the
left-hand side is $0$ and, since $(\binom di-1)^{d-i}t^d\neq0$ in
$\mathbb{Q}[t]$, we get $\varphi_i(c_i)=0$.
\end{proof}

These conditions are the natural starting point for a search for explicit
certificates over $\mathbb{Q}$: they fix the behaviour of each coefficient $c_j$
on $d-1$ explicit lines.

\subsection{Relation to Ghosh's proof}

Ghosh \cite{Gh25} announced a complete proof of the conjecture (Koszul homology
plus Brouwer degree arguments), revised in March 2026 and, as far as we could
see, not published in a journal; an August 2026 preprint still treats it as open.
If Ghosh's proof is correct, it implies Theorem~\ref{thm:all}, and what this note
contributes is an elementary and independent proof. Should it contain an error,
Theorem~\ref{thm:all} stands as an unconditional partial result extending
Theorem~\ref{thm:SS}.

\section{Computational verification}

Everything was verified with SageMath 10.7 in the exact setting of \cite{SS}
(Hasse derivatives, $f$ unnormalised except for $a_d=0$).

\begin{center}
\renewcommand{\arraystretch}{1.2}
\begin{tabular}{>{\raggedright\arraybackslash}p{0.34\textwidth}%
                >{\raggedright\arraybackslash}p{0.56\textwidth}}
\toprule
Statement & Verification\\
\midrule
Prop.~\ref{prop:lines}
  & all $284$ resultants $R_j(f_{i,t})$, $3\le d\le10$, all $i,j$, with $t$
    symbolic: no failures.\\
Theorem~\ref{thm:all}, by evaluation
  & $3\le d\le9$, all $i$.\\
Theorem~\ref{thm:all}, by Gröbner bases
  & a radical test independent of the witnesses (Rabinowitsch's trick,
    $1\in J_i+(1-yR_i)$): $d=3,4,5$, all $i$, with outcome ``not a member''.\\
Theorem~\ref{thm:all}, in Lean~4
  & all $d$ and all $i$, over any field of characteristic zero, with no
    \texttt{sorry}.\\
\bottomrule
\end{tabular}
\end{center}

\section{Caveats}

\begin{itemize}
\item We checked the statement of Theorem~\ref{thm:SS} in version~7 of \cite{SS}
  (February 2025). It is striking that the general case is not observed there; we
  do not rule out that it appears elsewhere, or that the interest of \cite{SS}
  lay in the real-root method rather than in the statement.
\item This is not a deep result: it is a lemma. Its main use is
  Corollary~\ref{cor:certificates}, which constrains membership certificates over
  $\mathbb{Q}$, and Lemma~\ref{lem:components}, which locates the difficulty
  precisely.
\end{itemize}

\begin{thebibliography}{9}

\bibitem{CA} E.~Casas--Alvero, \emph{Higher order polar germs},
  J.~Algebra \textbf{240} (2001), 326--337.

\bibitem{CLO} W.~Castryck, R.~Laterveer, M.~Ounaïes,
  \emph{Constraints on counterexamples to the Casas--Alvero conjecture and a
  verification in degree 12}, arXiv:1208.5404.

\bibitem{DdJ} J.~Draisma, J.~P.~de Jong, \emph{On the Casas--Alvero conjecture},
  Newsletter of the EMS \textbf{80} (2011), 29--33.

\bibitem{Gh24} S.~Ghosh, \emph{A finiteness result towards the Casas--Alvero
  conjecture}, arXiv:2402.18717.

\bibitem{Gh25} S.~Ghosh, \emph{Proof of the Casas--Alvero conjecture},
  arXiv:2501.09272 (v2, March 2026).

\bibitem{GLSV} H.-C.~Graf von Bothmer, O.~Labs, J.~Schicho,
  C.~van de Woestijne, \emph{The Casas--Alvero conjecture for infinitely many
  degrees}, J.~Algebra \textbf{316} (2007), 224--230.

\bibitem{SS} D.~Schaub, M.~Spivakovsky, \emph{A note on the Casas--Alvero
  conjecture}, arXiv:2312.08742 (v7, February 2025).

\end{thebibliography}

\end{document}
